\documentclass[crop,tikz]{standalone}
\usetikzlibrary{backgrounds}
\colorlet{blue}{cyan}
\tikzset{
  inverted/.style = {
    color=white,
    background rectangle/.style={fill},
    show background rectangle
  }
}

% Plots field lines of a bar magnet of length h using the expressions
% from Martín-Luna et.al., "On the Magnetic Field of a Finite
% Solenoid", IEEE TRANSACTIONS ON MAGNETICS, VOL. 59, NO. 4, APRIL
% 2023

% Note that the expressions are written in terms of the dimensionless
% ratio r/r_0. Thus, the x-axis shows the ratio x/r_0 and the y-axis
% shows the ratio y/r_0. I.e. the borders of the solenoid are at x=-1
% and x=1.

\usepackage{pgfplots}
\usepackage{luacode}

\pgfplotsset{
  compat=1.18,
  inverted/.style = {
    every axis legend/.append style={
      draw=white,
      fill=black,
      text=white
    }
  },
}

\begin{luacode*}
Ell = Ell or {}

-- 16-pt Gauss-Legendre nodes/weights on [-1,1], reduced to 8 positive nodes
Ell.xi = {
  0.9894009349916499,
  0.9445750230732326,
  0.8656312023878317,
  0.7554044083550030,
  0.6178762444026437,
  0.4580167776572274,
  0.2816035507792589,
  0.09501250983763744
}
Ell.wi = {
  0.027152459411754095,
  0.062253523938647893,
  0.095158511682492785,
  0.12462897125553387,
  0.14959598881657673,
  0.16915651939500254,
  0.18260341504492359,
  0.18945061045506850
}

Ell.mid  = math.pi/4   -- midpoint of [0,pi/2]
Ell.half = math.pi/4   -- half-length of [0,pi/2]
Ell.eps  = 1e-15

function Ell.clamp(x) if x < Ell.eps then return Ell.eps else return x end end

function Ell.K(k)
  local kk = k*k
  local sum = 0.0
  for i=1,#Ell.xi do
    local th1 = Ell.mid + Ell.half*Ell.xi[i]
    local th2 = Ell.mid - Ell.half*Ell.xi[i]
    local s1  = math.sin(th1); s1 = s1*s1
    local s2  = math.sin(th2); s2 = s2*s2
    local r1  = Ell.clamp(1.0 - kk*s1)
    local r2  = Ell.clamp(1.0 - kk*s2)
    sum = sum + Ell.wi[i]*(1.0/math.sqrt(r1) + 1.0/math.sqrt(r2))
  end
  return Ell.half*sum
end

function Ell.E(k)
  local kk = k*k
  local sum = 0.0
  for i=1,#Ell.xi do
    local th1 = Ell.mid + Ell.half*Ell.xi[i]
    local th2 = Ell.mid - Ell.half*Ell.xi[i]
    local s1  = math.sin(th1); s1 = s1*s1
    local s2  = math.sin(th2); s2 = s2*s2
    local r1  = Ell.clamp(1.0 - kk*s1)
    local r2  = Ell.clamp(1.0 - kk*s2)
    sum = sum + Ell.wi[i]*(math.sqrt(r1) + math.sqrt(r2))
  end
  return Ell.half*sum
end

function Ell.Pi(n,k)
  local kk = k*k
  local sum = 0.0
  for i=1,#Ell.xi do
    local th1 = Ell.mid + Ell.half*Ell.xi[i]
    local th2 = Ell.mid - Ell.half*Ell.xi[i]
    local s1  = math.sin(th1); s1 = s1*s1
    local s2  = math.sin(th2); s2 = s2*s2
    local r1  = Ell.clamp(1.0 - kk*s1)
    local r2  = Ell.clamp(1.0 - kk*s2)
    local d1  = Ell.clamp(1.0 - n*s1)
    local d2  = Ell.clamp(1.0 - n*s2)
    sum = sum + Ell.wi[i]*(1.0/(d1*math.sqrt(r1)) + 1.0/(d2*math.sqrt(r2)))
  end
  return Ell.half*sum
end
\end{luacode*}

\makeatletter
\pgfmathdeclarefunction{luaK}{1}{%
  \begingroup
  \pgfmathfloatparsenumber{#1}%
  \pgfmathfloattofixed{\pgfmathresult}%
  \edef\Karg{\pgfmathresult}%
  \pgfmathparse{%
    \directlua{
      local k = tonumber("\Karg")
      local fmt = string.char(37)..".16g"
      tex.sprint(string.format(fmt, Ell.K(k)))
    }%
  }%
  \pgfmath@smuggleone\pgfmathresult
  \endgroup
}
\pgfmathdeclarefunction{luaE}{1}{%
  \begingroup
  \pgfmathfloatparsenumber{#1}%
  \pgfmathfloattofixed{\pgfmathresult}%
  \edef\Earg{\pgfmathresult}%
  \pgfmathparse{%
    \directlua{
      local k = tonumber("\Earg")
      local fmt = string.char(37)..".16g"
      tex.sprint(string.format(fmt, Ell.E(k)))
    }%
  }%
  \pgfmath@smuggleone\pgfmathresult
  \endgroup
}
\pgfmathdeclarefunction{luaPi}{2}{%
  \begingroup
  \pgfmathfloatparsenumber{#1}%
  \pgfmathfloattofixed{\pgfmathresult}%
  \edef\Narg{\pgfmathresult}%
  \pgfmathfloatparsenumber{#2}%
  \pgfmathfloattofixed{\pgfmathresult}%
  \edef\Karg{\pgfmathresult}%
  \pgfmathparse{%
    \directlua{
      local n = tonumber("\Narg")
      local k = tonumber("\Karg")
      local fmt = string.char(37)..".16g"
      tex.sprint(string.format(fmt, Ell.Pi(n,k)))
    }%
  }%
  \pgfmath@smuggleone\pgfmathresult
  \endgroup
}
\makeatother

\begin{document}
\begin{tikzpicture}[inverted,inverted]
  \pgfmathsetmacro{\height}{2}; % ratio h/r_0
  \pgfmathsetmacro{\xmin}{-3};
  \pgfmathsetmacro{\xmax}{3};
  \pgfmathsetmacro{\ymin}{-3};
  \pgfmathsetmacro{\ymax}{3};
  \pgfmathsetmacro{\BO}{4*pi};
  \begin{axis}[inverted,
    xmin = {\xmin}, xmax = {\xmax},
    ymin = {\ymin}, ymax = {\ymax},
    axis equal image,
    view = {0}{90},
    hide axis,
    declare function={
      kk(\r,\z,\s) = sqrt(4*\r/((1 + \r)^2 + (\z + 0.5*\s*\height)^2)); % k_\pm
      kp(\r,\z) = kk(\r, \z, 1); % k_+
      km(\r,\z) = kk(\r, \z, -1); % k_-
      sigma(\r) = 4*\r/(1 + \r)^2;
      %
      Ipm(\r,\z,\k) = -4/sqrt(\r)*(luaE(\k)/\k - (1 - 0.5*\k^2)*luaK(\k)/\k); % I_\pm
      Jpm(\r,\z,\k,\s) = \k/sqrt(\r)*(\z + 0.5*\s*\height)*((1 - \r)/(1 + \r)*luaPi(sigma(\r),\k) + luaK(\k)); % I'_\pm
      Ip(\r,\z) = Ipm(\r,\z,kp(\r,\z)); % I_+
      Im(\r,\z) = Ipm(\r,\z,km(\r,\z)); % I_-
      Jp(\r,\z) = Jpm(\r,\z,kp(\r,\z),1); % I'_+
      Jm(\r,\z) = Jpm(\r,\z,km(\r,\z),-1); % I'_-
      %
      Bz0(\r,\z) = 0.5*\BO*((\z + 0.5*\height)/sqrt(1 + (\z + 0.5*\height)^2) - (\z - 0.5*\height)/sqrt(1 + (\z - 0.5*\height)^2)); % B_z(0,z)
      Br(\r,\z) = \r == 0 ? 0 : \BO/(4*pi)*(Im(\r,\z) - Ip(\r,\z)); % B_r(r,z)
      Brz(\r,\z) = \r == 0 ? Bz0(\r,\z) : \BO/(4*pi)*(Jp(\r,\z) - Jm(\r,\z)); % B_z(r,z)
      rr(\x,\y,\z) = sqrt(\x^2 + \y^2); % r(x,y,z)
      Bx(\x,\y,\z) = rr(\x,\y,\z) == 0 ? 0 : Br(rr(\x,\y,\z),\z)*cos(atan2(\y,\x)); % B_x(x,y,z)
      By(\x,\y,\z) = rr(\x,\y,\z) == 0 ? 0 : Br(rr(\x,\y,\z),\z)*sin(atan2(\y,\x)); % B_y(x,y,z)
      Bz(\x,\y,\z) = Brz(rr(\x,\y,\z),\z); % B_z(x,y,z)
      Ba(\x,\y,\z) = sqrt(Br(rr(\x,\y,\z),\z)^2 + Bz(\x,\y,\z)^2); % |\vec{B}(x,y,z)|
    },
    ]
    % magnetic vector field
    \addplot3[
      red,
      samples=51,
      samples y=51,
      quiver={
        u={Bx(x,0,y)},
        v={Bz(x,0,y)},
        scale arrows=0.04,
      },
    ] {0};
    % solenoid
    \draw[very thick] (-1,{-\height/2}) -- (-1,{\height/2});
    \draw[very thick] (1,{-\height/2}) -- (1,{\height/2});
  \end{axis}
\end{tikzpicture}
\end{document}
